\documentclass{article}
\usepackage{amsmath, amssymb, amsthm}
\usepackage{hyperref}
\usepackage{geometry}
\geometry{margin=1in}

\title{Erd\H{o}s Problem \#102}
\date{Last updated: 2025-08-31}
\author{Source: erdosproblems.com}

\begin{document}
\maketitle

\noindent\textbf{Status:} OPEN \quad \textbf{No prize}

\noindent\textbf{Tags:} geometry

\medskip

\noindent\textbf{Problem Statement:}

Let $c&#62;0$ and $h_c(n)$ be such that for any $n$ points in $\mathbb{R}^2$ such that there are $\geq cn^2$ lines each containing more than three points, there must be some line containing $h_c(n)$ many points. Estimate $h_c(n)$. Is it true that, for fixed $c&#62;0$, we have $h_c(n)\to \infty$?




A problem of Erd\H{o}s and Purdy. It is not even known if $h_c(n)\geq 5$ (see [101] ). It is easy to see that $h_c(n) \ll_c n^{1/2}$, and Erd\H{o}s at one point \cite{Er95} suggested that perhaps a similar lower bound $h_c(n)\gg_c n^{1/2}$ holds. Zach Hunter has pointed out that this is false, even replacing $&gt;3$ points on each line with $&gt;k$ points: consider the set of points in $\{1,\ldots,m\}^d$ where $n\approx m^d$. These intersect any line in $\ll_d n^{1/d}$ points, and have $\gg_d n^2$ many pairs of points each of which determine a line with at least $k$ points. This is a construction in $\mathbb{R}^d$, but a random projection into $\mathbb{R}^2$ preserves the relevant properties. This construction shows that $h_c(n) \ll n^{1/\log(1/c)}$.




References

[Er95] Erd\H{o}s, Paul, Some of my favourite problems in number theory, combinatorics, and geometry . Resenhas (1995), 165-186.


\bigskip
\noindent\small{Source: \url{https://www.erdosproblems.com/102}}
\end{document}
